\documentclass[9pt]{article}

\usepackage[margin=0.6in]{geometry}
\usepackage{amsmath,amssymb}
\usepackage{graphicx}
\usepackage{enumitem}
\usepackage{hyperref}
\usepackage{titlesec}
\usepackage{xcolor}

% compact layout
\setlength{\parindent}{0pt}
\setlength{\parskip}{1.5pt}
\renewcommand{\baselinestretch}{0.90}

\usepackage{enumitem}
\setlist{nosep,leftmargin=1.1em}

\usepackage{titlesec}
\titlespacing*{\section}{0pt}{3pt}{1.5pt}

\title{Project Plan: Certificate-Gated Defenses for Adversarial Control-Flow Injection in LLM Agents}
\author{}
\date{}

\begin{document}
\maketitle
\vspace{-2.0em}

\section*{Abstract}

Prompt and tool injection in large language model (LLM) agents enable adversarial content embedded in retrieved documents or tool outputs to induce unauthorized actions.
We formalize injection as an adversarial control-flow optimization problem in which an attacker perturbs untrusted input channels to maximize unsafe tool calls while preserving task plausibility.
We propose a certificate-gated authorization protocol that constrains actions via externally grounded certificates, reframing injection as constrained certificate forgery.
We analyze how this protocol strictly dominates tool allowlists, characterize adaptive attackers, and outline an experimental plan to evaluate robustness under adversarial budgets.

\section{Threat Model and Attack Formulation}

At step $t$, the agent observes $o_t=(s_t,r_t,u_t)$ consisting of internal state, retrieved content, and tool output.
The attacker perturbs untrusted channels; our evaluation focuses on retrieval injection:
\[
r_t' = r_t + \delta_t^{(r)}, \quad \delta\in\mathcal{C}.
\]
The formalism extends symmetrically to tool-output perturbations $u_t' = u_t + \delta_t^{(u)}$; experimental evaluation of this channel is left to future work.

Let $\pi_\theta(a_t\mid o_t')$ be the induced policy.
Let $\mathcal{B}$ denote unauthorized tool actions (e.g., exfiltration or policy-violating writes).

Attack success is defined as:
\[
\max_{\delta\in\mathcal{C}}
\mathbb{E}\Big[\sum_{t=1}^T \mathbf{1}\{a_t\in\mathcal{B}\}
-\lambda\ell_{\mathrm{task}}(a_{1:T})\Big].
\]

We restrict the attacker by a token budget $|\delta|\le B$ and at most $K$ injected sources per episode, and sweep $(B,K)$ in evaluation.

For analysis and estimation we define:
\[
\mathcal{L}_{\mathrm{bad}}(o)=\sum_{\tau\in\mathcal{T}_{\mathrm{bad}}}p_\theta(\tau\mid o),
\]
and empirically test correlation between $\Delta\mathcal{L}_{\mathrm{bad}}$ and realized attack rates.
Since computing the logit sum $p_\theta(\tau\mid o)$ requires model internals unavailable in all inference backends, we approximate $\mathcal{L}_{\mathrm{bad}}$ with a taint-score proxy based on n-gram overlap between action content and known payloads.

The verifier and external task specification are part of the trusted computing base and are not attacker-controlled.

\section{Certificate-Gated Authorization}

For experimental clarity, we assume a structured external task specification $\mathcal{G}$ (objective, inputs, permitted operations) derived deterministically from the user request.

For each privileged action the model emits $(a,\phi)$ with certificate:
\[
\phi=(g,E,C),
\]
where $g$ is a goal reference, $E$ cited evidence with provenance labels, and $C$ invoked constraints.

A deterministic verifier $V(a,\phi;\mathcal{G})$ accepts iff:

\begin{enumerate}[leftmargin=1.2em]
\item $\tau(a)$ satisfies capability constraints (schemas and tool permissions),
\item $g\in\Gamma(\mathcal{G})$ (externally grounded goals),
\item $C$ contains no tokens tainted by untrusted provenance,
\item untrusted content appears only inside quoted $E$.
\end{enumerate}

Taint is detected via n-gram overlap between $C$ and untrusted sources, augmented by cosine similarity of sentence embeddings; we note this provides only approximate paraphrase robustness.

This induces a constrained policy:
\[
\tilde{\pi}_\theta(a\mid o)\propto\sum_\phi \pi_\theta(a,\phi\mid o)\mathbf{1}\{V=1\}.
\]

\section{Strict Dominance Over Tool Allowlists}

Let $\pi^{\text{allow}}$ be an allowlist baseline enforcing only $\tau(a)\in\mathcal{T}_{\text{allow}}$.

Certificate gating enforces:
\[
\mathcal{F}_{\text{cert}}=\{(a,\phi):\tau(a)\in\mathcal{T}_{\text{allow}}\land V(a,\phi;\mathcal{G})=1\}.
\]

Therefore:
\[
\Pr_{\tilde{\pi}_\theta}[a\in\mathcal{B}\cap\mathcal{T}_{\text{allow}}]
\le
\Pr_{\pi^{\text{allow}}}[a\in\mathcal{B}\cap\mathcal{T}_{\text{allow}}],
\]
since certificate gating rejects a subset of actions permitted by allowlists but lacking valid authorization.

We quantify this gap via:
\[
\widehat{\Delta}_{\text{auth}}=\mathbb{E}[\mathbf{1}\{\tau(a)\in\mathcal{T}_{\text{allow}}\land V=0\}],
\]
measuring how often certificate gating blocks attacks that allowlists alone would permit.

\section{Adaptive Attacker as Certificate Forgery}

Under the defense, attackers must induce:
\[
\max_{\delta\in\mathcal{C}}
\mathbb{E}\Big[\sum_t \mathbf{1}\{a_t\in\mathcal{B}\land V(a_t,\phi_t;\mathcal{G})=1\}\Big].
\]

We evaluate three adaptive strategies:

\begin{enumerate}[leftmargin=1.2em]
\item \textbf{Goal laundering:} attempting to redefine objectives to satisfy $g\in\Gamma(\mathcal{G})$.
\item \textbf{Evidence laundering:} embedding instructions inside cited documents.
\item \textbf{Policy mimicry:} imitating constraint language to bypass $C$ checks.
\end{enumerate}

\section{Estimators and Mechanism-Level Metrics}

From execution traces we compute:
\[
\widehat{R}_{\mathrm{bad}}=\Pr[a\in\mathcal{B}],\quad
\widehat{R}_{\mathrm{forge}}=\Pr[a\in\mathcal{B}\land V=1],
\]
with bootstrap confidence intervals over episodes.

We correlate $\widehat{R}_{\mathrm{forge}}$ with $\Delta\mathcal{L}_{\mathrm{bad}}$ to test whether certificate gating flattens the attacker's optimization landscape.

We additionally log verifier rejection modes and action divergence from clean trajectories to diagnose failure mechanisms.

\section{Experimental System and Evaluation}

Primary evaluation focuses on single-agent systems under indirect prompt injection via retrieval.
Secondary experiments extend to planner-executor agents.

Baselines include input quoting, provenance labeling, static allowlists, and their combinations.

Primary metrics:
$\widehat{R}_{\mathrm{bad}}$, $\widehat{R}_{\mathrm{forge}}$, task success, and $\widehat{\Delta}_{\text{auth}}$.
We report security-utility frontiers by varying a taint-threshold parameter in the verifier.

\section{Project Plan}

\textbf{Weeks 1-2: Foundations.}
Build agent harness with full trace logging.
Implement indirect injection and tool injection attacks.
Establish clean baselines.

\textbf{Weeks 3-4: Adversarial Suite.}
Construct injected document and tool-output corpus.
Implement adaptive attack strategies under budgets $(B,K)$.
Validate attack effectiveness without defenses.

\textbf{Weeks 5-6: Certificate-Gated Defense.}
Implement certificate schema, taint tracking, and deterministic verifier.
Integrate constrained policy $\tilde{\pi}_\theta$.

\textbf{Weeks 7-8: Comparative Evaluation.}
Benchmark certificate gating against baselines.
Measure $\widehat{R}_{\mathrm{bad}}$, $\widehat{R}_{\mathrm{forge}}$, and utility.

\textbf{Weeks 9-10: Mechanism Analysis.}
Analyze verifier rejection modes and attribution to injected tokens.
Run ablations on verifier components.


\textbf{Weeks 11-12: Write-up and Release.}
Finalize report, package code and execution traces, and prepare demo.

\section{Expected Outcomes and Limitations}

The project will deliver:

\begin{itemize}[leftmargin=1.2em]
\item a formal adversarial formulation of control-flow injection,
\item a certificate-gated authorization protocol for LLM agents,
\item empirical robustness results against adaptive attackers,
\item security-utility tradeoff curves,
\item code and execution traces for reproducibility.
\end{itemize}

We note that models may learn to satisfy verifier constraints over time; robustness under joint training is left to future work.

\section{Broader Impact}

As LLM agents enter scientific, financial, and healthcare workflows, preventing untrusted content from directly triggering actions is critical.
Certificate-gated authorization provides a practical, model-agnostic defense that can be deployed without retraining, offering a concrete step toward safer agentic systems.

\end{document}
